\section{Attack Surfaces}
\label{sec:attack}

Three attack surfaces correspond to the three data sources.
Table~\ref{tab:attacks} summarises the attacks, their magnitude, and
their detectability under the \bisect{} audit chain.

\begin{table}[h]
\centering
\caption{Input data manipulation attack surfaces.
         Detection probability under \bisect{} audit chain.
         $\Delta D$: estimated maximum Democratic seat shift, national.}
\label{tab:attacks}
\begin{tabular}{llccc}
\toprule
Attack & Data source & $\Delta D$ & Detection & Method \\
\midrule
Population inflation/deflation & P.L.\ 94-171 & $\leq 0.4$ & $\approx 1.00$ & Download hash \\
Shapefile boundary modification & TIGER/Line & $\leq 0.4$ & $\approx 0.98$ & Adjacency hash \\
ACS demographic alteration & ACS 5-year & $\leq 0.1$ & $\approx 1.00$ & Download hash \\
\midrule
Combined (empirical, saturated) & All three & $\leq 0.4$ & $\approx 0.98$ & \\
\bottomrule
\end{tabular}
\end{table}

All three attacks require the adversary to control the data pipeline
between the Census Bureau's servers and the \bisect{} computation.
This typically means the adversary must either (a) compromise the
redistricting authority's internal data storage, (b) intercept network
traffic during the download phase (the TLS limitation discussed in
Section~\ref{sec:tls}), or (c) substitute files after download but
before the hash is computed.

Option (c) is the most feasible in practice but is prevented by computing
the SHA-256 hash immediately upon download and storing it in the manifest.
Options (a) and (b) require infrastructure-level compromise that is
beyond the scope of redistricting-specific gaming.
